\chapter{Instrument and Sequence Examples}
\label{app:examples}

This appendix gives worked examples of instruments and sequences: a lead, a bass, a violin with vibrato, a kick drum, a
hat and a snare, and a plucked note, then a one-bar pattern that plays several of them. Each one lists everything a
program writes to make the sound, and what the sound does. They are starting points. Change a waveform, an envelope
step or a modulation depth, and listen.

\section{Reading the Examples}

\subsection{Instrument Records}

An instrument is described by the record of Table~\ref{tab:instrument}. The examples give each record as the values of
its fields. \texttt{SAMPLE\_START}, \texttt{SAMPLE\_LENGTH} and \texttt{SPIT\_ADDRESS} depend on where the program puts
the waveform and the pitch table in audio memory, so only the length is given. Every waveform in the examples is a loop
that starts at byte 0 and covers the whole sample, so \texttt{LOOP\_START} is 0 and \texttt{LOOP\_ENABLE} is set (word
6 is \texttt{0x8000}).

\subsection{Pitch Tables}

The pitch table of an instrument holds the twelve steps of one octave (Section~\ref{sec:pitch-table}). A table for a
one-cycle waveform of $N$ bytes has entries $N/64$ times those of Table~\ref{tab:pitch-table}.
Table~\ref{tab:example-spit} gives the entry for note 0 of each table the examples use. Instruments with the same
length of waveform share a table.

\begin{table}[H]
	\centering
	\begin{tabular}{lll}
		\toprule
		\textbf{Waveform}  & \textbf{Entry 0 (A0)} & \textbf{Used by}                  \\
		\midrule
		64 bytes           & \texttt{0x00000E19}   & Lead, kick                        \\
		256 bytes          & \texttt{0x00003864}   & Bass, violin, pluck               \\
		4096 bytes         & \texttt{0x0003863E}   & Hat, snare                        \\
		64 bytes, $\div 5$ & \texttt{0x000002D2}   & Vibrato source (5.5~Hz at note 0) \\
		\bottomrule
	\end{tabular}
	\caption{Pitch Tables Used in the Examples}
	\label{tab:example-spit}
\end{table}

\subsection{Envelope Times}

An envelope takes one step every 2.0~ms (Chapter~\ref{ch:envelope}), so the time a phase takes is the number of steps
times 2.0~ms. Attack takes $\lceil 4095 / \texttt{ATTACK} \rceil$ steps. Decay takes $\lceil (4095 - S) /
	\texttt{DECAY} \rceil$ steps, where $S$ is the 12-bit sustain level, and release takes $\lceil S / \texttt{RELEASE}
	\rceil$ steps. An \texttt{ATTACK} of 0 is instant. A \texttt{DECAY} of 0 goes to the sustain level at once.

\subsection{Conditions}

The levels and times quoted are for a voice that plays alone, at a volume of 255 and a pan of 128, with
\texttt{MASTER\_VOL} at 255 and \texttt{AUDIO\_EN} set. A note is started by writing \texttt{INSTR} and then
\texttt{COMMAND} (Section~\ref{sec:startup}). Where a sequence is used, the tick is 400 samples, 12.5~ms
(\texttt{SEQ\_PERIOD} = 400).

\section{A Square Lead}

A plain tone with a soft start. The waveform is a square of 64 bytes, $+90$ in the first half and $-90$ in the second.

\begin{table}[H]
	\centering
	\begin{tabular}{llp{6cm}}
		\toprule
		\textbf{Field}                     & \textbf{Value} & \textbf{Meaning}                                     \\
		\midrule
		\texttt{SAMPLE\_LENGTH}            & 64             & One cycle                                            \\
		\texttt{ATTACK}, \texttt{DECAY}    & 128, 8         & Word 4 is \texttt{0x8008}                            \\
		\texttt{SUSTAIN}, \texttt{RELEASE} & 190, 24        & Word 5 is \texttt{0xBE18}; sustain level 3051 (75\%) \\
		\bottomrule
	\end{tabular}
	\caption{Square Lead: Instrument}
	\label{tab:ex-lead}
\end{table}

The voice needs \texttt{CONTROL} = \texttt{0x0000} and \texttt{VOLUME\_PAN} = \texttt{0xFF80}. To play A4, write the
instrument number to \texttt{INSTR} and \texttt{0x8030} (note 48) to \texttt{COMMAND}.

\textbf{Result.} The note sounds at 440.0~Hz. The attack takes 32 steps, 64~ms, and the level then decays for about
262~ms to the sustain level. After a stop the release takes 256~ms to reach silence, and the level is below 1\% of the
peak after 252~ms.

\section{A Filtered Bass}

A sawtooth of 256 bytes, rising from $-105$ to $+105$, through the low-pass filter. The sawtooth is rich in harmonics
and the filter takes the top off them.

\begin{table}[H]
	\centering
	\begin{tabular}{llp{6cm}}
		\toprule
		\textbf{Field}                     & \textbf{Value} & \textbf{Meaning}                                     \\
		\midrule
		\texttt{SAMPLE\_LENGTH}            & 256            & One cycle                                            \\
		\texttt{ATTACK}, \texttt{DECAY}    & 0, 8           & Word 4 is \texttt{0x0008}; instant attack            \\
		\texttt{SUSTAIN}, \texttt{RELEASE} & 140, 40        & Word 5 is \texttt{0x8C28}; sustain level 2248 (55\%) \\
		\bottomrule
	\end{tabular}
	\caption{Filtered Bass: Instrument}
	\label{tab:ex-bass}
\end{table}

\begin{table}[H]
	\centering
	\begin{tabular}{lll}
		\toprule
		\textbf{Register}          & \textbf{Value}  & \textbf{Meaning}                                       \\
		\midrule
		\texttt{CONTROL}           & \texttt{0x2000} & \texttt{FILTER\_EN}: the voice goes through the filter \\
		\texttt{VOLUME\_PAN}       & \texttt{0xFF80} & Volume 255, centre                                     \\
		\texttt{FILTER\_CONFIG}    & 0               & Low-pass                                               \\
		\texttt{FILTER\_CUTOFF}    & 16              & 318~Hz (Table~\ref{tab:filter-cutoff})                 \\
		\texttt{FILTER\_RESONANCE} & 64              & A gentle peak at the cutoff                            \\
		\texttt{FILTER\_VOL\_PAN}  & \texttt{0xFF80} & Filter volume 255, centre                              \\
		\bottomrule
	\end{tabular}
	\caption{Filtered Bass: Registers}
	\label{tab:ex-bass-regs}
\end{table}

\textbf{Result.} Note 12 sounds A1 at 55.0~Hz. The fundamental is unchanged by the filter. The sixteenth harmonic,
880~Hz, nearly three octaves above the cutoff, is 16.6~dB weaker than in the same voice with \texttt{FILTER\_EN}
clear. A higher cutoff brightens the sound. Writing the cutoff register while the note plays gives the familiar sweep,
in steps at the low end (Section~\ref{sec:filter-cutoff}).

\section{A Violin with Vibrato}

A bowed tone has a slow attack and a vibrato that comes in after the note has started. It needs two voices: one plays
the tone, and a second, in modulation mode, supplies the vibrato (Chapter~\ref{ch:modulation}). The tone is a sawtooth
of 256 bytes with an amplitude of 100. The vibrato source is a sine of 64 bytes with an amplitude of 120.

\begin{table}[H]
	\centering
	\begin{tabular}{llll}
		\toprule
		\textbf{Field}                     & \textbf{Tone (instrument 0)} & \textbf{Vibrato source (instrument 1)} \\
		\midrule
		\texttt{SAMPLE\_LENGTH}            & 256                          & 64                                     \\
		Pitch table                        & 256 bytes                    & 64 bytes, $\div 5$                     \\
		\texttt{ATTACK}, \texttt{DECAY}    & 30, 4                        & 8, 0                                   \\
		Word 4                             & \texttt{0x1E04}              & \texttt{0x0800}                        \\
		\texttt{SUSTAIN}, \texttt{RELEASE} & 200, 20                      & 255, 10                                \\
		Word 5                             & \texttt{0xC814}              & \texttt{0xFF0A}                        \\
		\bottomrule
	\end{tabular}
	\caption{Violin: Instruments}
	\label{tab:ex-violin}
\end{table}

\begin{table}[H]
	\centering
	\begin{tabular}{llll}
		\toprule
		\textbf{Register}    & \textbf{Voice 0 (tone)} & \textbf{Voice 5 (vibrato)} & \textbf{Meaning}                      \\
		\midrule
		\texttt{CONTROL}     & \texttt{0x0000}         & \texttt{0x4000}            & Voice 5 is in modulation mode         \\
		\texttt{VOLUME\_PAN} & \texttt{0xFF80}         & \texttt{0xFF80}            & Volume 255, centre                    \\
		\texttt{MOD\_ROUTE}  & ---                     & \texttt{0x2000}            & Pitch, voice 0, shift 0, not inverted \\
		\bottomrule
	\end{tabular}
	\caption{Violin: Registers}
	\label{tab:ex-violin-regs}
\end{table}

Start the vibrato source with \texttt{INSTR} = 1 and \texttt{COMMAND} = \texttt{0x8000} on voice 5. Start the tone with
\texttt{INSTR} = 0 and \texttt{COMMAND} = \texttt{0x8030} on voice 0. A fingered phrase is a run of \texttt{NOTE}
instructions on voice 0; the vibrato source keeps running underneath.

\textbf{Result.} The tone sounds A4 at 440~Hz. The tone's attack takes 137 steps, 274~ms. The vibrato source runs at
5.5~Hz, since note 0 of its table is built for that rate. Its own attack is slow, 512 steps or 1.0~s, so the vibrato
grows with it: in the first 200~ms the pitch moves by about $\pm 2.5$ cents, and from the first second it moves by
about $\pm 14$ cents.

The depth is set by the shift in \texttt{MOD\_ROUTE}, and each step of the shift doubles it. Pitch modulation is added
in raw units (Section~\ref{sec:mod-pitch-depth}), so the same setting gives half the depth in cents an octave higher. A
tone with a shorter waveform needs a smaller shift for the same depth in cents.

\section{A Kick Drum}

A kick is a low sine whose pitch falls quickly. The sine has 64 bytes and an amplitude of 115. The note starts high and
a second \texttt{NOTE} glides it down, which is what \texttt{PORTA\_EN} is for (Section~\ref{sec:note-start}).

\begin{table}[H]
	\centering
	\begin{tabular}{llp{6cm}}
		\toprule
		\textbf{Field}                     & \textbf{Value} & \textbf{Meaning}                                 \\
		\midrule
		\texttt{SAMPLE\_LENGTH}            & 64             & One cycle                                        \\
		\texttt{ATTACK}, \texttt{DECAY}    & 0, 60          & Word 4 is \texttt{0x003C}                        \\
		\texttt{SUSTAIN}, \texttt{RELEASE} & 0, 20          & Word 5 is \texttt{0x0014}; the level decays to 0 \\
		\bottomrule
	\end{tabular}
	\caption{Kick: Instrument}
	\label{tab:ex-kick}
\end{table}

The voice needs \texttt{CONTROL} with \texttt{PORTA\_EN} set (\texttt{0x1000}, or \texttt{0x9000} with the sequencer
on), \texttt{VOLUME\_PAN} = \texttt{0xB080} and \texttt{PORTA\_RATE} = 1120. One hit is this sequence:

\begin{table}[H]
	\centering
	\small
	\begin{tabular}{llp{6.4cm}r}
		\toprule
		\textbf{Word}   & \textbf{Instruction} & \textbf{Effect}                                   & \textbf{Tick} \\
		\midrule
		\texttt{0x0024} & \texttt{NOTE 36}     & Starts A3, 220~Hz                                 & 0             \\
		\texttt{0x4001} & \texttt{WAIT 1}      & One tick                                          &               \\
		\texttt{0x000C} & \texttt{NOTE 12}     & Same instrument, voice active: glides to A1       & 1             \\
		\texttt{0x4013} & \texttt{WAIT 19}     & Lets the sound die away                           &               \\
		\texttt{0x2000} & \texttt{STOP}        & Releases the envelope, so the next hit retriggers & 20            \\
		\bottomrule
	\end{tabular}
	\caption{One Kick Hit}
	\label{tab:ex-kick-seq}
\end{table}

The first \texttt{NOTE} finds the voice idle, so it starts a note. The second finds the voice active on the same
instrument with \texttt{PORTA\_EN} set, so it glides. The \texttt{STOP} is needed because a \texttt{NOTE} on a voice
that is still active would glide again and not start a new hit; with the envelope released and finished, the next
\texttt{NOTE} retriggers.

\textbf{Result.} The pitch falls from 220~Hz to 55~Hz in 38~ms. The time of a glide is about the distance in step
units divided by \texttt{PORTA\_RATE}, times 2.0~ms; here $21{,}654 / 1120 \times 2.0 = 38.7$~ms. The level falls below 1\%
of its peak after 136~ms.

\section{A Hat and a Snare}

Noise needs no pitch, only a sample that is not periodic. Both instruments use 4096 bytes of noise with an amplitude of
100, loop the whole sample, and share one voice by changing the instrument between hits. The hat and the snare use
different noise, and the note number only chooses which bytes are read. Use a pitch table for 4096 bytes
(Table~\ref{tab:example-spit}).

\begin{table}[H]
	\centering
	\begin{tabular}{lll}
		\toprule
		\textbf{Field}                     & \textbf{Hat (instrument 1)} & \textbf{Snare (instrument 2)} \\
		\midrule
		\texttt{ATTACK}, \texttt{DECAY}    & 0, 255 (\texttt{0x00FF})    & 0, 40 (\texttt{0x0028})       \\
		\texttt{SUSTAIN}, \texttt{RELEASE} & 0, 255 (\texttt{0x00FF})    & 0, 60 (\texttt{0x003C})       \\
		Note used                          & 100                         & 62                            \\
		\bottomrule
	\end{tabular}
	\caption{Hat and Snare: Instruments}
	\label{tab:ex-hat}
\end{table}

A hit is \texttt{INST 1} (\texttt{0x6001}) then \texttt{NOTE 100} (\texttt{0x0064}) for the hat, or \texttt{INST 2}
(\texttt{0x6002}) then \texttt{NOTE 62} (\texttt{0x003E}) for the snare.

\textbf{Result.} Both sounds start at full level. The hat falls below 2\% of its peak after 32~ms, and the snare after
202~ms.

\section{A Plucked Note}

A short, bright note for an arpeggio. The waveform is a triangle of 256 bytes that starts at zero and rises, with an
amplitude of 105.

\begin{table}[H]
	\centering
	\begin{tabular}{llp{6cm}}
		\toprule
		\textbf{Field}                     & \textbf{Value} & \textbf{Meaning}          \\
		\midrule
		\texttt{SAMPLE\_LENGTH}            & 256            & One cycle                 \\
		\texttt{ATTACK}, \texttt{DECAY}    & 0, 80          & Word 4 is \texttt{0x0050} \\
		\texttt{SUSTAIN}, \texttt{RELEASE} & 0, 60          & Word 5 is \texttt{0x003C} \\
		\bottomrule
	\end{tabular}
	\caption{Pluck: Instrument}
	\label{tab:ex-pluck}
\end{table}

\textbf{Result.} Note 51 sounds C5 at 523.3~Hz. The level reaches silence 104~ms after the start (the decay takes 52
steps), and is below 1\% of the peak after 102~ms. A run of these notes, each longer than 110~ms, starts
each one from silence, so a retrigger makes no step in the output (Chapter~\ref{ch:envelope}).

\section{A One-Bar Pattern}

The last example plays three of the instruments together as a bar of sixteen steps and repeats it. A step is a
sixteenth note: ten ticks, 125.2~ms, which makes four steps a beat at 119.9 beats a minute. A bar is 160 ticks,
2002.5~ms. The voices are:

\begin{table}[H]
	\centering
	\begin{tabular}{llll}
		\toprule
		\textbf{Voice} & \textbf{Instruments}                    & \textbf{\texttt{CONTROL}} & \textbf{\texttt{VOLUME\_PAN}} \\
		\midrule
		0              & Kick (instrument 0)                     & \texttt{0x9000}           & \texttt{0xB080}               \\
		1              & Hat (1) and snare (2)                   & \texttt{0x8000}           & \texttt{0x5A80}               \\
		2              & Bass (3, as in Table~\ref{tab:ex-bass}) & \texttt{0x8000}           & \texttt{0xC880}               \\
		\bottomrule
	\end{tabular}
	\caption{One-Bar Pattern: Voices}
	\label{tab:ex-pattern-voices}
\end{table}

\noindent Voice 0 also has \texttt{PORTA\_RATE} = 1120. The bass voice does not use the filter, so its \texttt{FILTER\_EN} is clear. Write \texttt{SEQ\_PC} of each voice with the address of its sequence, and write
\texttt{SEQ\_PERIOD} = 400 last, so that all three sequences start on the first tick (Section~\ref{sec:startup}).

\subsection{The Kick Sequence}

The kick plays on steps 0, 8 and 11. After each hit comes a wait that brings the sequence to the next hit.

\begin{table}[H]
	\centering
	\small
	\begin{tabular}{p{6.7cm}p{4.3cm}r}
		\toprule
		\textbf{Words}                                     & \textbf{Instructions}                                                                                  & \textbf{Tick} \\
		\midrule
		\texttt{0x6000}                                    & \texttt{INST 0}                                                                                        & 0             \\
		\texttt{0x0024 0x4001 0x000C 0x4013 0x2000 0x403C} & \texttt{NOTE 36}, \texttt{WAIT 1}, \texttt{NOTE 12}, \texttt{WAIT 19}, \texttt{STOP}, \texttt{WAIT 60} & 0 (step 0)    \\
		\texttt{0x0024 0x4001 0x000C 0x4013 0x2000 0x400A} & the same hit, then \texttt{WAIT 10}                                                                    & 80 (step 8)   \\
		\texttt{0x0024 0x4001 0x000C 0x4013 0x2000 0x401E} & the same hit, then \texttt{WAIT 30}                                                                    & 110 (step 11) \\
		\texttt{0xBFED}                                    & \texttt{GOTO} $-19$                                                                                    & 160           \\
		\bottomrule
	\end{tabular}
	\caption{Kick Sequence}
	\label{tab:ex-kick-pattern}
\end{table}

\noindent The \texttt{GOTO} jumps back 19 words, to the first \texttt{NOTE}, not to the \texttt{INST}, which need
only be selected once (Section~\ref{sec:seq-instructions}).

\subsection{The Hat and Snare Sequence}

The hat plays on steps 2, 6, 10 and 14 and the snare on steps 4 and 12. Each hit is preceded by the wait from the hit
before it, so the sequence begins with a wait.

\begin{table}[H]
	\centering
	\small
	\begin{tabular}{p{5.4cm}p{5.4cm}r}
		\toprule
		\textbf{Words}                & \textbf{Instructions}                                & \textbf{Step}            \\
		\midrule
		\texttt{0x4014 0x6001 0x0064} & \texttt{WAIT 20}, \texttt{INST 1}, \texttt{NOTE 100} & 2 (hat)                  \\
		\texttt{0x4014 0x6002 0x003E} & \texttt{WAIT 20}, \texttt{INST 2}, \texttt{NOTE 62}  & 4 (snare)                \\
		\texttt{0x4014 0x6001 0x0064} & \texttt{WAIT 20}, \texttt{INST 1}, \texttt{NOTE 100} & 6 (hat)                  \\
		\texttt{0x4028 0x6001 0x0064} & \texttt{WAIT 40}, \texttt{INST 1}, \texttt{NOTE 100} & 10 (hat)                 \\
		\texttt{0x4014 0x6002 0x003E} & \texttt{WAIT 20}, \texttt{INST 2}, \texttt{NOTE 62}  & 12 (snare)               \\
		\texttt{0x4014 0x6001 0x0064} & \texttt{WAIT 20}, \texttt{INST 1}, \texttt{NOTE 100} & 14 (hat)                 \\
		\texttt{0x4014 0xBFEC}        & \texttt{WAIT 20}, \texttt{GOTO} $-20$                & 16, the start of the bar \\
		\bottomrule
	\end{tabular}
	\caption{Hat and Snare Sequence}
	\label{tab:ex-hat-pattern}
\end{table}

\subsection{The Bass Sequence}

The bass plays on every second step, root then octave, as notes 12 and 24. Each note is held for 19 ticks and released,
and one tick is left before the next, so that every note starts from a released envelope.

\begin{table}[H]
	\centering
	\small
	\begin{tabular}{p{6.7cm}p{4.3cm}r}
		\toprule
		\textbf{Words}                       & \textbf{Instructions}                                              & \textbf{Tick} \\
		\midrule
		\texttt{0x6003}                      & \texttt{INST 3}                                                    & 0             \\
		\texttt{0x000C 0x4013 0x2000 0x4001} & \texttt{NOTE 12}, \texttt{WAIT 19}, \texttt{STOP}, \texttt{WAIT 1} & 0             \\
		\texttt{0x0018 0x4013 0x2000 0x4001} & \texttt{NOTE 24}, \texttt{WAIT 19}, \texttt{STOP}, \texttt{WAIT 1} & 20            \\
		\multicolumn{3}{l}{The last two lines are repeated three more times, and make eight notes in all (ticks 0 to 140).}       \\
		\texttt{0xBFDF}                      & \texttt{GOTO} $-33$                                                & 160           \\
		\bottomrule
	\end{tabular}
	\caption{Bass Sequence}
	\label{tab:ex-bass-pattern}
\end{table}

\textbf{Result.} The kick sounds at the start of the bar, at step 8 and at step 11, and the pattern repeats every
160 ticks with no drift. The three sequences complete every tick without a missed deadline. The sequences loop for
ever, so they never set \texttt{SEQ\_DONE}; to finish a song, end each voice with \texttt{END 0} (Chapter~\ref{ch:sequencer}).

\section{Summary}

\begin{table}[H]
	\centering
	\small
	\begin{tabular}{llllp{4.6cm}}
		\toprule
		\textbf{Example} & \textbf{Waveform} & \textbf{ADSR}   & \textbf{Note} & \textbf{Measured}                        \\
		\midrule
		Lead             & Square, 64        & 128, 8, 190, 24 & A4, 440.0~Hz  & Attack 64~ms, release 252~ms to 1\%      \\
		Bass             & Sawtooth, 256     & 0, 8, 140, 40   & A1, 55.0~Hz   & 16th harmonic $-16.6$~dB with the filter \\
		Violin           & Sawtooth, 256     & 30, 4, 200, 20  & A4, 440~Hz    & Vibrato 5.5~Hz, $\pm 14$ cents after 1~s \\
		Kick             & Sine, 64          & 0, 60, 0, 20    & A3 to A1      & 220 to 55~Hz in 38~ms; 1\% after 136~ms  \\
		Hat              & Noise, 4096       & 0, 255, 0, 255  & 100           & 2\% after 32~ms                          \\
		Snare            & Noise, 4096       & 0, 40, 0, 60    & 62            & 2\% after 202~ms                         \\
		Pluck            & Triangle, 256     & 0, 80, 0, 60    & C5, 523.3~Hz  & 1\% after 102~ms                         \\
		\bottomrule
	\end{tabular}
	\caption{Summary of the Examples}
	\label{tab:ex-summary}
\end{table}
